\scanpage{48}
% archive sheet number: 47
\[
0\unsure{\text{---}|\to}P'\mathbin{\text{---}}P\mathbin{\text{---}}P''\mathbin{\text{---}}0\qquad Q
\]
\deleted{\unsure{$nP''=0$|$nP''$}} a) $\,\unsure{n|{}_n}Q=Q$\quad b) $\,{}_nP\to P''$ \unsure{épimorphisme|épijectif}

\margin{\deleted{c) la suite $0\to P'\to P\to P''\to0$ est splittée \deleted{\unsure{ou}} \uline{ou} elle est \deleted{$\operatorname{Hom}(\unsure{P|P'}\otimes Q,G)\to\operatorname{Hom}(P'\otimes Q,G)$} \uline{loc. splittée} et $\underline{\operatorname{Hom}}(P'\otimes Q,G)=0$}}

Toute biextension \inserted{$E$} de $P\times Q$ \deleted{\unsure{définit|déf.}} par $G$ définit

a) Une biextension \inserted{$E'$} de $P'\times Q$ par $G$ d'où un accouplement
\[
{}_nP'\times{}_nQ\xrightarrow{E'(n)}G
\]

b) Une biextension de $\,{}_nP\times Q$ par $G$, i.e. un accouplement
\[
{}_nP\times{}_nQ\xrightarrow{\varphi}G
\]

\deleted{\unsure{avec la|avec le|Avec les}} Ces deux $(E',\varphi)$ satisfont la condition que
\[
\varphi\mid({}_nP'\times{}_nQ)=E'(n)
\]

Un hom de biextensions $E_1\to E_2$ définit un hom $(E'_1,\varphi_1)\unsure{\text{---}|\to}(E'_2,\varphi_2)$, i.e. \inserted{(déf)} un hom de biextensions $u:E'_1\to E'_2$, tel que $\varphi\,\unsure{\sim|\text{---}}$

\uline{Je dis que le foncteur} $E\mapsto(E',\varphi)$ est \uline{une équiv. de catégories}.

\deleted{\deleted{a) \unsure{pl. fid.}} \deleted{b) pl. fid.} \deleted{i.e. a) $\operatorname{Hom}(P\otimes Q,G)\xrightarrow{\sim}\operatorname{Hom}(P'\otimes Q,G)$ b) $E_1\simeq E_2$ \unsure{ssi|ss} $(E'_1,\varphi_1)\simeq(E'_2,\varphi_{\unsure{1|2}})$}

\uline{Pl. \unsure{fidél.|fidélité|fidèle}} Signifie aussi que $E\mathrel{\shortmid\!-}E|(P',Q)$ est pl. fid, on commencera \deleted{\unsure{donc|par}} par la suite exacte, que
\[
\unsure{\circledast|(*)}\quad\operatorname{Hom}(P''\otimes Q,G)=0,\qquad
(**)\quad\operatorname{Hom}(P\otimes Q,G)\to\operatorname{Hom}(P'\otimes Q,G)\ \uuline{\text{surj.}}
\]
\deleted{Ce qui est évidemment \unsure{ainsi|aussi} $(***)$ Si \illegible{1} tel que $(E',\varphi)\unsure{\simeq|=}\unsure{1|0}$ alors $\unsure{E'|E}\simeq\unsure{E_1|1}$.}
Si $E$ est tel que $E'\simeq1,\varphi=0$, alors $E\simeq1$.
La première relation résulte de a) et b) qui impliquent $P''\otimes Q=0$, la deuxième résulte de c) [dans la deuxième alternative de c), les deux mbres de $(**)$ sont nuls !].

\deleted{\uline{\unsure{Ess. surjectif|Ess. surject.}}}. Pour $(***)$, on note que si $E'=0$, alors par la suite exacte $E$ est déterminé par un élément de $\operatorname{Ext}^1(P''\mathbin{\overset{L}{\otimes}}Q,G)$, or ceci équivaut à un élément de $\operatorname{Hom}(P''\otimes_{\unsure{1|-}} Q,G)$, on devine que c'est ce qui définit $\varphi$, on gagne.}

\uline{Corollaire} 1 Supposons, pour un $\mathbb P$ de nbs premiers que $(\varinjlim_{\substack{n\\\text{\unsure{à support|support} }\mathbb P}}{}_nP)\to P''$ surjectif, et que $p\in\mathbb P\Rightarrow pQ=Q$.

Alors on a une équivalence $E\mapsto(E',\varphi)$, où $E$ est une
